% Vierfeldertafel – bank/vierfeldertafel/ (zusammenbau v0.4, Heft)
\einheitenkopf[t3]{Vierfeldertafel}
\setcounter{aufgabe}{16}

% Anlauf (ohne Prüfkennung)
\begin{aufgabe}{Die Tafel füllen – Anlauf}
\begin{teile}
\teil In einem Sportverein gilt: „$40\,\%$ der Mitglieder spielen Tennis.“ Was ist diese Angabe? \\ \kreuz{Rand} \\ \kreuz{Feld} \\ \kreuz{bedingter Anteil}
\teil In einer Stadt haben $40\,\%$ der Haushalte ein Zeitungsabo ($Z$), $30\,\%$ sind Seniorenhaushalte ($S$), und $18\,\%$ sind Seniorenhaushalte mit Zeitungsabo. Vervollständige die Vierfeldertafel (in Prozent).

\vierfeldertafel{S}{Z}{18,,40,,,,30,,100}
\teil In einem Chor sind $60\,\%$ der Mitglieder Frauen ($F$); $35\,\%$ aller Mitglieder singen Sopran ($S$); von den Frauen singt die Hälfte Sopran. Vervollständige die Vierfeldertafel (in Prozent).

\vierfeldertafel{S}{F}{,,60,,,,35,,100}
\end{teile}
\end{aufgabe}

\begin{aufgabe}{Die Tafel füllen – Prüfungsaufgaben}
\begin{teile}
\teil Von den Beschäftigten eines Klinikums sind insgesamt $64\,\%$ weiblich ($W$); $35\,\%$ arbeiten in der Pflege ($K$), und $22\,\%$ der Beschäftigten sind weiblich und arbeiten in der Pflege. Vervollständige die Vierfeldertafel (in Prozent). \hfill (Abitur 2023 GK)

\vierfeldertafel{K}{W}{22,,64,,,,35,,100}
\teil Von den Kunden eines Fitnessstudios haben $70\,\%$ einen Jahresvertrag ($J$), $30\,\%$ trainieren vor allem abends ($A$), und $8\,\%$ haben keinen Jahresvertrag und trainieren vor allem abends. Vervollständige die Vierfeldertafel (in Prozent). \hfill (Abitur 2024 GK)

\vierfeldertafel{A}{J}{,,70,8,,,30,,100}
\teil In einem Versandlager sind $20\,\%$ der Pakete Expresspakete ($E$), $12\,\%$ aller Pakete kommen verspätet an ($V$), und $25\,\%$ der Expresspakete kommen verspätet an. Vervollständige die Vierfeldertafel (in Prozent). \hfill (Abitur 2022 GK)

\vierfeldertafel{V}{E}{,,20,,,,12,,100}
\teil Bei einem Leihfahrrad ist mit $8\,\%$ Wahrscheinlichkeit die Bremse defekt ($B$), mit $5\,\%$ das Licht ($L$), und mit $89\,\%$ ist weder die Bremse noch das Licht defekt. Vervollständige die Vierfeldertafel (in Prozent). \hfill (Abitur 2018 GK)

\vierfeldertafel{L}{B}{,,8,,89,,5,,100}
\end{teile}
\end{aufgabe}

% Anlauf (ohne Prüfkennung)
\begin{aufgabe}{Aus der Tafel rechnen – Anlauf}
\begin{teile}
\teil $A$: Ein Kunde kauft Brot, $B$: er kauft Brötchen. Gesucht ist die Wahrscheinlichkeit für „Brot oder Brötchen“. Welcher Ausdruck passt? \\ \kreuz{$1 - P(\overline{A} \cap \overline{B})$} \\ \kreuz{$P(A \cap \overline{B}) + P(\overline{A} \cap B)$} \\ \kreuz{$P(\overline{A} \cap \overline{B})$}
\teil Von $300$ Gästen eines Hotels sind $120$ Familien ($F$); $180$ Gäste buchen Halbpension ($H$), davon $90$ Familien. Wie viele Gäste sind keine Familie und buchen keine Halbpension? \leerfeld
\teil Von den Besuchern eines Museums haben $45\,\%$ eine Jahreskarte ($J$), $30\,\%$ nutzen den Audioguide ($G$), und $12\,\%$ tun beides. Welcher Anteil der Besucher hat eine Jahreskarte oder nutzt den Audioguide? \leerfeld[\%]
\end{teile}
\end{aufgabe}

\begin{aufgabe}{Aus der Tafel rechnen – Prüfungsaufgaben}
\begin{teile}
\teil Zu den Sporttests eines Landes traten $9\,640$ Personen an, davon $1\,875$ im Alter von mindestens $40$ Jahren; $7\,210$ bestanden, davon $6\,030$ jünger als $40$ Jahre. Wie viele Teilnehmer waren jünger als $40$ Jahre und haben nicht bestanden? \hfill (Abitur 2019 GK) \\ \leerfeld
\teil Von den Studierenden einer Hochschule wohnen $58\,\%$ in der Stadt ($W$), $30\,\%$ studieren ein technisches Fach ($T$), und $21\,\%$ tun beides. Mit welcher Wahrscheinlichkeit wohnt eine zufällig ausgewählte Person in der Stadt oder studiert ein technisches Fach? \hfill (Abitur 2023 GK) \\ \leerfeld[\%]
\teil Bei einem Sportfest sind $70\,\%$ der Teilnehmer Kinder, die übrigen Erwachsene. $14\,\%$ aller Teilnehmer sind Erwachsene mit Medaille; von den Kindern hat die Hälfte eine Medaille. Welcher Anteil aller Teilnehmer hat keine Medaille? \hfill (Abitur 2020 GK) \\ \leerfeld[\%]
\end{teile}
\end{aufgabe}
